\section{Conclusions}\label{conclusions}

\subsection{Summary of the Study}\label{summary-of-the-study}

This study examined how residential property values in Metro Cebu could be estimated more consistently using a data-driven workflow with geospatial feature engineering. The problem began with the gap between fast-moving market prices and slower administrative or appraisal benchmarks. To address that problem, the study assembled a cleaned analytical base table from open-market listings, bank ROPA inventory, floor-price references, BIR zonal values, and geospatial layers derived from Google Maps and OpenStreetMap. The final deployed prediction scope was limited to open-market residential property in six LGUs: Cebu City, Mandaue City, Lapu-Lapu City, Talisay City, Minglanilla, and Consolacion.

The study followed the CRISP-DM process and compared three modeling approaches: OLS hedonic regression, Random Forest, and XGBoost. The workflow also constructed a Metro Cebu Residential Accessibility Index (MCRAI) using category-specific radii and gravity-style decay, then paired that index with polycentric distance features, structural property variables, and spatial context variables. The held-out evaluation results showed that the baseline Random Forest model produced the strongest overall performance on the retained deployment criteria, with $R^2 = 0.807$, MAE = 4.95 million PHP, and RMSE = 27.45 million PHP on the test set. The thesis therefore concluded that a machine-learning workflow grounded in local urban structure can produce a usable decision-support tool for Metro Cebu open-market residential valuation, even if prediction error remains substantial in some property types and locations.

\subsection{Answers to the Research Questions}\label{answers-to-the-research-questions}

Research Question 1 asked, ``What value drivers significantly influence property prices in Metro Cebu?'' The study found that both locational and property-type variables were important, but the strongest recurring drivers were geospatial rather than purely structural. In the retained tree models, property-type indicators such as single detached, condominium, and vacant lot occupied the top positions. Distance to Consolacion was the strongest locational distance variable in both models, ranking consistently in the top four features. Distance to Cebu Business Park also appeared in the top ten for both models, while longitude ranked among the top features in XGBoost. The conclusion was that property prices in Metro Cebu were shaped by a combination of polycentric accessibility and housing submarket segmentation, not by one uniform citywide gradient.

Research Question 2 asked, ``Which modeling technique---\textbf{Hedonic Regression}, \textbf{Random Forest}, or \textbf{XGBoost}---produces the most accurate valuation out of sample (lowest MAPE)?'' The final comparison made that question more nuanced than the original wording suggested. The confirmation-tuned Random Forest achieved the lowest MAPE at 53.00\%, but it did so with materially weaker overall fit than the deployed baseline Random Forest. On the retained multi-metric selection rule, the baseline Random Forest produced the strongest overall held-out result on the 299-row test set, with $R^2 = 0.807$, MAE = 4.95 million PHP, and RMSE = 27.45 million PHP. OLS performed much worse as a predictive model, while both XGBoost variants remained below the retained baseline on the overall deployment criteria. The conclusion was therefore two-part: the lowest MAPE was delivered by the tuned Random Forest, but the strongest overall out-of-sample valuation model in the final workflow was the baseline Random Forest, which was retained as the production model for the Streamlit application.

Research Question 3 asked, ``Do \textbf{geospatial features}---proximity to economic nodes, amenity density, and spatial autocorrelation---significantly improve model performance compared to structural-only models?'' The study answered this question substantially, though not in the form of a standalone structural-only ablation table. The final evidence showed that geospatial variables were central to the fitted models rather than marginal additions. Polycentric distance features dominated SHAP attribution, longitude remained highly influential, and the MCRAI variables contributed to the econometric screening stage that produced the final positive-category composite. These findings supported the conclusion that geospatial features materially improved the explanatory and practical value of the modeling workflow. At the same time, the final manuscript did not retain a separate structural-only benchmark table, so the exact incremental uplift from geospatial features alone was not isolated as a single chapter-level number.

Research Question 4 asked, ``How large is the ``Valuation Gap'' between the model's data-driven predictions and traditional BIR Zonal Values?'' The study answered this question operationally more than statistically. The workflow defined \emph{valuation gap} as the difference between the model-facing price level and the residential BIR zonal benchmark, and this field was retained as a diagnostic output in the analytical base table and spatial outputs. This made the gap observable at the property and map level. However, the final results chapters did not report one metro-wide summary figure for the valuation gap distribution. The conclusion was therefore narrower: the thesis established a defensible way to measure and visualize the gap, but it did not reduce that gap to a single final statistic for the whole study area.

\subsection{Methodological Contribution}\label{methodological-contribution}

The methodological contribution of the thesis rests on the way established valuation ideas were adapted to the specific conditions of Metro Cebu. First, the study scoped the production model to open-market residential property rather than forcing a single predictive target across open-market, bank ROPA, and floor-price observations. That decision kept the deployed output aligned with a market-value use case and avoided mixing fundamentally different pricing regimes into the final prediction surface.

Second, the thesis built a polycentric distance feature set instead of relying on one central-business-district distance measure. This mattered because the final findings showed that Metro Cebu behaved more like a corridor with several active nodes than like a purely monocentric market. Third, the thesis introduced the two-stage MCRAI workflow. Stage 1 used OLS signs and coefficients to evaluate the individual accessibility categories. Stage 2 retained only the positive-coefficient household-access categories in the final composite and treated the negative categories as diagnostic evidence of spatial sorting rather than as simple disamenities. Together, these decisions produced a locally grounded valuation workflow that was more specific to Metro Cebu than a generic hedonic template.

\subsection{Practical Contribution}\label{practical-contribution}

The practical contribution of the thesis is the decision-support layer that came out of the modeling work. The main deliverable is a Streamlit price surface for open-market residential property in Metro Cebu. This output does not replace formal appraisal judgment, but it gives brokers, appraisers, lenders, and researchers a reproducible spatial estimate of residential market price under a clearly defined market segment.

The thesis also produced QGIS-ready layers that organized the model outputs spatially. These layers support map-based review of predicted price levels and related diagnostics across the study area. Within that workflow, the \emph{valuation gap} field remains useful because it allows users to compare modeled or observed market-oriented prices against BIR zonal benchmarks without treating the benchmark itself as the prediction target. In that sense, the thesis contributes not only a model, but also a practical way to inspect where official benchmarks and market-facing estimates may diverge.

\subsection{Closing Statement}\label{closing-statement}

The thesis concluded that data science methods can support residential valuation in Metro Cebu when they are applied with clear market definitions, local geospatial logic, and disciplined treatment of value concepts. The strongest results did not come from model complexity alone. They came from matching the workflow to the structure of the local market: an open-market target, a polycentric urban form, and an accessibility framework designed for household valuation rather than for liquidation or administrative pricing. Under those conditions, the study produced a defensible decision-support tool for Metro Cebu residential valuation.

